% ============================================================================
% Fichier  : partie6/P06_C19_retex_amelioration.tex
% Titre    : Retour d'Expérience et Amélioration Continue
% Auteur   : MINKA MI NGUIDJOÏ Thierry Emmanuel
% ============================================================================
% Sources : ISO/IEC 27035 (phase d'enseignements) ; culture juste. Op. MAYO.
% ============================================================================

\chapter{Retour d'Expérience et Amélioration Continue}
\label{chap:retex}
\chaptermark{Retour d'expérience}

\epigraph{%
    « Un incident dont on ne tire pas de leçon\\
    est un incident que l'on paiera deux fois. »%
}{Principe du cours}

\niveauChapitre{Fondamental}{Partie~V.}

\begin{boxobjectifs}
\begin{itemize}
    \item Conduire un retour d'expérience utile et sans procès.
    \item Distinguer causes immédiates et causes profondes.
    \item Transformer les leçons en actions suivies.
    \item Comprendre la culture juste.
    \item Boucler sur la supervision et la détection.
\end{itemize}
\end{boxobjectifs}

\bigskip

% ============================================================================
\section*{Question directrice : qu'a appris LiZbiZ, et comment éviter que cela se reproduise~?}
% ============================================================================

Une fois le service rétabli, la tentation est de tourner la page. C'est
l'erreur~: la phase d'enseignements de l'ISO 27035 est celle qui fait
progresser. Ce chapitre conduit le RETEX et le transforme en actions.

\medskip
\textbf{Plan du chapitre.}
\begin{enumerate}
    \item Conduire un RETEX (section~\ref{sec:c19-conduire}).
    \item Causes immédiates et profondes (section~\ref{sec:c19-causes}).
    \item Des leçons aux actions (section~\ref{sec:c19-actions}).
    \item La culture juste (section~\ref{sec:c19-culture}).
\end{enumerate}

\separateur

% ============================================================================
\section{Conduire un RETEX}
\label{sec:c19-conduire}
% ============================================================================

\begin{boxprocedure}
\textbf{Conduire un retour d'expérience}
\begin{enumerate}
    \item Réunir les acteurs dans les 72~heures, à partir de la main
          courante.
    \item Reconstituer la chronologie des faits (pas des responsabilités).
    \item Identifier ce qui a bien et mal fonctionné.
    \item Remonter aux causes profondes, au-delà du symptôme.
    \item Décider des actions, avec responsable et échéance.
\end{enumerate}
\end{boxprocedure}

% ============================================================================
\section{Causes Immédiates et Profondes}
\label{sec:c19-causes}
% ============================================================================

\begin{tableaucours}{L{4.4cm}L{7.8cm}}{%
    \tableauHeader{Niveau} &
    \tableauHeader{Exemple (GAROUA)}
}
Cause immédiate & Un compte de service à mot de passe faible compromis \\
Cause intermédiaire & Pas de supervision, journaux non lus \\
Cause profonde & Absence de fonction de détection et de gouvernance de la
sécurité \\
\end{tableaucours}
\vspace{-0.8em}
\begin{center}{\small\itshape Tableau~19.1. Remonter aux causes profondes}\end{center}

\begin{boxalert}
\textbf{Erreur fréquente.} S'arrêter à la cause immédiate (« le mot de passe
était faible ») et sanctionner.

\smallskip
\textbf{Règle~:} corriger la seule cause immédiate laisse la cause profonde
intacte~; l'incident reviendra sous une autre forme. On traite les trois
niveaux.
\end{boxalert}

% ============================================================================
\section{Des Leçons aux Actions}
\label{sec:c19-actions}
% ============================================================================

\begin{boxscenario}{Opération MAYO. Le plan d'action}
\textbf{Constatation.} Le RETEX de LiZbiZ débouche sur~: mettre en place une
supervision (SOC hybride), journaliser \texttt{FS-01}, durcir les comptes de
service, sauvegarder hors ligne, écrire des playbooks, former le personnel.
Chaque action a un responsable et une échéance.

\medskip
\textbf{Analyse} (confiance élevée). Les actions traitent les trois niveaux
de causes. Le RETEX n'est utile que si ses actions sont suivies~: LiZbiZ les
inscrit dans un plan revu chaque trimestre.
\end{boxscenario}

% ============================================================================
\section{La Culture Juste}
\label{sec:c19-culture}
% ============================================================================

\begin{boxdef}
\termecle{Culture juste}~: culture où l'on signale les erreurs et les
quasi-incidents sans crainte de sanction injuste. Elle distingue l'erreur
honnête (que l'on analyse) de la faute délibérée (que l'on sanctionne). Sans
elle, les incidents se cachent et ne s'apprennent pas.
\end{boxdef}

\begin{boxmetier}{RSSI}
\textbf{Protéger ceux qui signalent.} Le RSSI fait du RETEX un espace sans
procès~: le technicien qui a redémarré \texttt{FS-01} de bonne foi n'est pas
sanctionné, il devient l'exemple qui fonde la consigne de première
intervention. On apprend des erreurs, on ne les enterre pas.
\end{boxmetier}

\separateur

\begin{boxresume}
\textbf{Ce qu'il faut retenir de ce chapitre~:}
\begin{itemize}
    \item \textbf{Le RETEX} se conduit à partir de la main courante, sans
          procès.
    \item \textbf{Trois niveaux de causes}~: immédiate, intermédiaire,
          profonde.
    \item \textbf{On traite les trois}, sinon l'incident revient.
    \item \textbf{Des actions suivies}, avec responsable et échéance.
    \item \textbf{La culture juste} fait remonter ce qu'il faut apprendre.
\end{itemize}
\end{boxresume}

\begin{boxinfo}
\textbf{Points de contrôle de l'élève.} Je sais~: conduire un RETEX~;
distinguer les niveaux de causes~; transformer les leçons en actions~;
expliquer la culture juste~; boucler sur la supervision.
\end{boxinfo}

\bigskip

\begin{boxexo}[title={\ding{45}\quad Exercice~19.1. Délai \niveaubase}]
Dans quel délai conduire le RETEX~? À partir de quel document~?
\end{boxexo}

\begin{boxexo}[title={\ding{45}\quad Exercice~19.2. Causes \niveaubase}]
Distinguez cause immédiate et cause profonde dans GAROUA.
\end{boxexo}

\begin{boxexo}[title={\ding{45}\quad Exercice~19.3. Culture juste \niveaubase}]
Qu'est-ce que la culture juste~?
\end{boxexo}

\begin{boxexo}[title={\ding{45}\quad Exercice~19.4. Symptôme \niveaumoyen}]
Pourquoi corriger la seule cause immédiate ne suffit-il pas~?
\end{boxexo}

\begin{boxexo}[title={\ding{45}\quad Exercice~19.5. Actions \niveaumoyen}]
Reliez chaque action du plan de MAYO à un niveau de cause.
\end{boxexo}

\begin{boxexo}[title={\ding{45}\quad Exercice~19.6. Sans procès \niveaumoyen}]
Pourquoi le RETEX ne doit-il pas être un procès~?
\end{boxexo}

\begin{boxexo}[title={\ding{45}\quad Exercice~19.7. Le technicien \niveaumoyen}]
Comment traiter le cas du technicien qui a redémarré \texttt{FS-01}~?
\end{boxexo}

\begin{boxexo}[title={\ding{45}\quad Exercice~19.8. Suivi \niveaumoyen}]
Comment s'assurer que les actions du RETEX sont réellement suivies~?
\end{boxexo}

\begin{boxexo}[title={\ding{45}\quad Exercice~19.9. Compte rendu de RETEX \niveauexpert}]
Rédigez le compte rendu de RETEX de l'opération GAROUA.
\end{boxexo}

\begin{boxexo}[title={\ding{45}\quad Exercice~19.10. Plan d'amélioration \niveauexpert}]
Construisez le plan d'amélioration trimestriel de LiZbiZ issu du RETEX.
\end{boxexo}

\bigskip

\begin{boxinfo}
\textbf{Pour aller plus loin.}
\begin{itemize}
    \item ISO/IEC 27035~\cite{iso27035}.
\end{itemize}
\end{boxinfo}

\begin{boxinfo}
\textbf{Transition.} Reste à mesurer tout cela~: le
Chapitre~\ref{chap:mesurer} clôt le cours sur la mesure et la recherche.
\end{boxinfo}
